The rotating calipers algorithm~\citep{toussaint1983solving} computes the MBR
of a convex polygon in $O(n)$ time after the convex hull is available.
Since the convex hull is computed in $O(n \log n)$ (K.3, Section~3),
the total MBR computation is $O(n \log n)$.

\subsection{Algorithm}

\begin{algorithm}
\caption{Rotating Calipers for Minimum Bounding Rectangle}
\label{alg:rotating_calipers}
\begin{algorithmic}[1]
\Function{MinBoundingRect}{$H$: convex hull vertices}
  \State Find the four extremal vertices: top, right, bottom, left
  \State Initialise four calipers along cardinal directions
  \State $A_{\min} \gets \infty$; $\text{MBR} \gets \text{None}$
  \For{$i = 0$ to $|H|-1$} \Comment{rotate through each hull edge}
    \State Edge vector $\mathbf{e} \gets H[i+1] - H[i]$
    \State Rotate calipers so one caliper is flush with edge $\mathbf{e}$
    \State Compute $\ell, w$ from caliper positions (projections onto $\mathbf{e}$
      and its perpendicular)
    \If{$\ell \cdot w < A_{\min}$}
      \State $A_{\min} \gets \ell \cdot w$; $\text{MBR} \gets (\ell, w)$
    \EndIf
  \EndFor
  \State \Return MBR
\EndFunction
\end{algorithmic}
\end{algorithm}

At each step, the calipers advance to the next hull edge by rotating by the
minimum angle needed to align with that edge.
The total rotation is $2\pi$ over all $n$ hull edges, giving $O(n)$ steps.

\subsection{Example: Rotated Rectangle}

Consider a $1 \times 3$ rectangle rotated 45 degrees.
Its vertices are (in approximate coordinates):
$(0,0)$, $(0.71, -0.71)$, $(2.83, 1.41)$, $(2.12, 2.12)$.

AABB: $\min x = 0$, $\max x = 2.83$, $\min y = -0.71$, $\max y = 2.12$.
AABB dimensions: $2.83 \times 2.83$; LW$_{\text{AABB}} = 1.0$.

MBR: rotating calipers finds the rectangle oriented at 45 degrees with
dimensions $3 \times 1$; LW$_{\text{MBR}} = 3.0$.

AABB underestimates LW from $3.0$ to $1.0$ (a 100\% error) in this example.
For realistic districts oriented at modest angles (15--30 degrees),
the AABB overestimate is typically 5--41\%.

\subsection{Implementation}

Implemented in \texttt{crates/bisect-analysis/src/compactness.rs}.
The convex hull is computed first (Andrew's monotone chain, K.3 Section~3),
then rotating calipers is applied to the hull vertices.
Numerical precision: coordinates in projected metres (EPSG:5070);
LW reported to 3 decimal places.

The production compactness surface exposes
\texttt{length\_width\_ratio()} as the MBR-based metric used by
\texttt{all\_metrics()}.  The crate also exposes
\texttt{axis\_aligned\_length\_width\_ratio()} as a diagnostic helper for this
paper's AABB comparison.  A unit test verifies the central distinction: a
rotated $3{:}1$ rectangle retains MBR LW $=3$, while its AABB LW collapses to
$1$ at a $45^\circ$ orientation.
